% ======================================================================
% PREÁMBULO: Tesina de Licenciatura - Configuración Unificada
% ======================================================================

% --------------------------------------------------
% 1. Codificación, Fuentes e Idioma
% --------------------------------------------------
\usepackage[utf8]{inputenc}
\usepackage[T1]{fontenc}
\usepackage[spanish,es-tabla]{babel}

% Tipografía formal Times New Roman estándar
\usepackage{mathptmx}

% --------------------------------------------------
% 1.1. Color Institucional y Fondos de Portada (DEPENDENCIAS DE PORTADA)
% --------------------------------------------------
\usepackage[dvipsnames]{xcolor}
\definecolor{AzulBUAP}{RGB}{0, 50, 110} % Azul oficial BUAP[cite: 4]
\usepackage{eso-pic}                     % Permite colocar marcas de agua y líneas de fondo[cite: 4]

% --------------------------------------------------
% 2. Bibliografía APA 7 (BibLaTeX + Biber)
% --------------------------------------------------
\usepackage{csquotes}
\usepackage[
  backend=biber,
  style=apa,
  language=spanish
]{biblatex}
\addbibresource{referencias.bib}

% --------------------------------------------------
% 3. Geometría de Página
% --------------------------------------------------
\usepackage[
  letterpaper,
  margin=2.54cm,
  headheight=15pt,
  headsep=12pt
]{geometry}

% --------------------------------------------------
% 4. Formato de Párrafo e Interlineado
% --------------------------------------------------
\usepackage{setspace}
\doublespacing                         % Interlineado doble APA 7
\setlength{\parindent}{1.27cm}         % Sangría de primera línea uniforme
\setlength{\parskip}{0pt}
\setcounter{secnumdepth}{4}

% --------------------------------------------------
% 5. Formato de Títulos (APA 7)
% --------------------------------------------------
\usepackage{titlesec}
\titleformat{\chapter}[display]
  {\normalfont\huge\bfseries\centering}
  {\chaptertitlename\ \thechapter}
  {10pt}
  {\Huge}
\titlespacing*{\chapter}{0pt}{10pt}{25pt}

\titleformat{\section}
  {\normalfont\Large\bfseries}
  {\thesection}
  {1em}
  {}
\titlespacing*{\section}{0pt}{15pt}{8pt}

\titleformat{\subsection}
  {\normalfont\large\bfseries\itshape}
  {\thesubsection}
  {1em}
  {}

% --------------------------------------------------
% 6. Control de Encabezados y Foliación Uniforme
% --------------------------------------------------
\usepackage{emptypage}
\usepackage{fancyhdr}

\pagestyle{fancy}
\fancyhf{}
\renewcommand{\headrulewidth}{0pt}
\rhead{\thepage} % Número siempre fijo en la esquina superior derecha

\fancypagestyle{plain}{%
  \fancyhf{}%
  \renewcommand{\headrulewidth}{0pt}%
  \rhead{\thepage}%
}

% --------------------------------------------------
% 7. Tablas, Flotantes y Figuras
% --------------------------------------------------
\usepackage{booktabs}
\usepackage{longtable}
\usepackage{multirow}
\usepackage{graphicx}
\usepackage{float}
\usepackage{array}       
\usepackage{amsmath,amssymb,amsthm}
\newtheorem{theorem}{Teorema}[chapter]

\usepackage{caption}
\captionsetup{
  justification=raggedright,
  singlelinecheck=false,
  labelfont=bf,
  font=small
}
\usepackage{subcaption}

% --------------------------------------------------
% 8. Hipervínculos Compatibles con APA 7
% --------------------------------------------------
\usepackage[hidelinks]{hyperref}

% Permite que el enlace funcione sin romper la estructura de los paréntesis
\DeclareFieldFormat{citehyperref}{%
  \iffieldundef{postnote}
    {\bibhyperref{#1}}
    {\bibhyperref{#1}}}